\documentclass[11pt,a4paper]{article}

\usepackage[utf8]{inputenc}
\usepackage[main=italian,provide=*]{babel}
\usepackage{amsmath}
\usepackage{amssymb}
\usepackage{mathtools}
\usepackage{geometry}
\usepackage{booktabs}
\usepackage{hyperref}
\usepackage{seqsplit}
\usepackage{graphicx}

\geometry{margin=2.7cm}

\title{Il notebook \texttt{\seqsplit{quantum\_simulation\_trimero\_anello\_trotter.ipynb}} spiegato\\
\large Fisica, matematica e numeri, passo per passo, senza salti logici}
\author{Note per la discussione col relatore --- Tesi triennale, Samuele Schiavi\\ Relatore: Prof.\ Alessandro Chiesa}
\date{}

\newcommand{\ket}[1]{\lvert #1 \rangle}
\newcommand{\bra}[1]{\langle #1 \rvert}

\begin{document}
\maketitle
\tableofcontents

\bigskip
\noindent Questo documento accompagna, cella per cella, il notebook
\texttt{\seqsplit{quantum\_simulation\_trimero\_anello\_trotter.ipynb}}:
per ogni sezione del notebook, spiega \emph{cosa} viene calcolato,
\emph{perché} (la fisica e la matematica dietro), e \emph{cosa
significano} i numeri che escono. È pensato per essere portato al
colloquio col relatore come traccia puntuale, indipendente dagli altri
documenti di teoria (che vengono comunque citati per chi vuole la
derivazione completa di ogni singolo passaggio algebrico).

\section{Il problema e l'Hamiltoniana (cella di setup)}

Il notebook implementa la simulazione dell'evoluzione temporale reale
$e^{-iHt}$ per il trimero anello isoscele ($N=3$ spin-$1/2$), con
\begin{align}
H &= b\sum_{i=1}^3 Z_i \;+\; H_{ex} \;+\; H_{DM}, \label{eq:H-full}\\[3pt]
H_{ex} &= J\,\vec\sigma_1\!\cdot\!\vec\sigma_2 \;+\; J'\big(\vec\sigma_2\!\cdot\!\vec\sigma_3+\vec\sigma_3\!\cdot\!\vec\sigma_1\big), \notag\\[3pt]
H_{DM} &= D_{12}(X_1Z_2{-}Z_1X_2) \;+\; D_{23}(X_2Z_3{-}Z_2X_3) \;+\; D_{31}(X_3Z_1{-}Z_3X_1), \notag
\end{align}
in \textbf{convenzione di Pauli diretta} (nessun fattore $1/2$ o $1/4$:
gli operatori sono direttamente $X,Y,Z$, non $\sigma/2$), con base $J$ sui
siti $(1,2)$ e lati $J'$ su $(2,3),(3,1)$ (topologia isoscele), e Opzione
B del termine DM: $D_{12}=rJ$, $D_{23}=D_{31}=rJ'$ (nel codice il
parametro $r$ è chiamato \texttt{D}). Questa è \emph{l'estensione a
$N=3$} di un compito già completato per il dimero ($N=2$): l'obiettivo
del notebook è mostrare che, sebbene $H$ non si possa esponenziare in
blocco senza diagonalizzare una matrice $8\times8$, l'evoluzione si può
comunque implementare con un circuito quantistico, tramite decomposizione
di Suzuki--Trotter, e quantificare l'errore che questa approssimazione
introduce.

Il notebook importa due moduli:
\begin{itemize}
\item \texttt{\seqsplit{trimer\_ring\_exact.py}} -- il benchmark esatto già
usato per tutta la teoria e il VQE del trimero: costruisce $H$ come
\texttt{SparsePauliOp} e la diagonalizza esattamente (\texttt{eigh}),
serve qui come "verità di riferimento" contro cui misurare l'errore di
Trotter;
\item \texttt{\seqsplit{trotter\_trimero\_anello.py}} -- il modulo scritto per
questo compito, che costruisce il circuito Trotter e lo confronta con la
matrice esatta.
\end{itemize}
Il fatto che entrambi condividano la stessa costruzione
dell'Hamiltoniana (il secondo modulo importa letteralmente le funzioni
del primo, non ne ridefinisce una copia) è una scelta metodologica
precisa: elimina la possibilità che una discrepanza numerica sia dovuta a
una convenzione di segno o di indicizzazione diversa tra i due pezzi di
codice, invece che a un vero errore di approssimazione fisica.

\section{Self-test del modulo (Sezione 1 del notebook)}
\label{sec:self-test}

Prima di produrre qualunque risultato fisico, il notebook esegue quattro
verifiche indipendenti. Questo ordine non è casuale: se uno di questi
test fallisse, qualunque numero prodotto dopo sarebbe privo di senso, non
un errore di Trotter genuino ma un bug di implementazione (convenzione di
segno, fattore mancante, qubit scambiati).

\subsection{Self-test 1: coerenza dell'operatore $S_z^{tot}$}

Il notebook verifica che l'operatore $S_z^{tot}=\sum_i Z_i$, costruito
direttamente in \texttt{\seqsplit{trotter\_trimero\_anello.py}}, coincida con
quello ottenibile a partire dal benchmark esatto. Il trucco usato è
semplice ma rigoroso: l'Hamiltoniana a scambio nullo, $H(J{=}0,J'{=}0,b)$,
è per costruzione $b\sum_iZ_i=b\,S_z^{tot}$ (Eq.~\eqref{eq:H-full} con
$J=J'=0$); quindi la differenza $H(0,0,b_2)-H(0,0,b_1)$, divisa per
$b_2-b_1$, deve restituire esattamente $S_z^{tot}$, indipendentemente da
quale dei due moduli l'abbia costruito. Il risultato del notebook,
$\|S_z^{tot}(\text{qui})-S_z^{tot}(\text{da trimer\_hamiltonian})\|=0$,
conferma che non c'è disallineamento fra le due costruzioni.

\subsection{Self-test 2: il circuito Qiskit reale contro la matrice indipendente}
\label{subsec:selftest2}

Questo è il test più importante dei quattro, perché confronta
\textbf{il circuito che verrà davvero eseguito} (non solo una sua
rappresentazione a matrice teorica) con una costruzione a matrice
indipendente. Per cinque set di parametri casuali $(J,J',b,D,t)$, il
notebook:
\begin{enumerate}
\item costruisce il circuito Qiskit per un singolo passo di Trotter
(\texttt{trotter\_circuit(...,\,N{=}1)}) ed estrae la sua unitaria
$8\times8$ esatta tramite \texttt{Operator(qc).data} (Qiskit calcola
letteralmente il prodotto di tutte le matrici dei gate, nell'ordine in
cui compaiono nel circuito);
\item costruisce, \emph{separatamente e senza passare dal circuito}, la
stessa unitaria a partire dalla definizione matematica: esponenziale di
matrice (\texttt{scipy.linalg.expm}) di ciascun blocco di legame,
moltiplicati fra loro nell'ordine previsto dalla teoria;
\item confronta le due matrici con la norma dell'operatore.
\end{enumerate}
L'errore trovato è dell'ordine di $10^{-15}$ per tutti e cinque i casi --
precisione macchina, non un residuo di approssimazione fisica. Questo
significa che il circuito Qiskit implementa \emph{esattamente} la
costruzione matematica descritta nella Sezione~\ref{sec:trotter-error},
gate per gate: se ci fosse un errore di segno, di fattore (es.\ $\theta$
invece di $2\theta$ in un gate di rotazione), o di indicizzazione dei
qubit, l'errore qui misurato sarebbe dell'ordine dell'unità, non
$10^{-15}$.

\subsection{Self-test 3: convergenza $N\to\infty$ verso l'evoluzione esatta}
\label{subsec:selftest3}

Fissato un punto $(J,J',b,D,t)$, il notebook calcola l'infedeltà
$1-|\langle\psi_{\rm esatto}(t)|\psi_{\rm Trotter}(t)\rangle|^2$ per
$N=5,10,20,40,80$ passi, partendo da $\ket{000}$. I numeri prodotti:
\begin{center}
\begin{tabular}{@{}ccc@{}}
\toprule
$N$ & infedeltà & rapporto rispetto al precedente \\
\midrule
5 & $6.091\times10^{-3}$ & --- \\
10 & $1.129\times10^{-3}$ & $5.39$ \\
20 & $2.614\times10^{-4}$ & $4.32$ \\
40 & $6.387\times10^{-5}$ & $4.09$ \\
80 & $1.585\times10^{-5}$ & $4.03$ \\
\bottomrule
\end{tabular}
\end{center}
Il rapporto tende a $4$ (non a $2$) al crescere di $N$: questo è
\emph{atteso} per un metodo del prim'ordine, e la ragione è spiegata per
esteso nella Sezione~\ref{sec:trotter-error}. In breve: l'errore
sull'\emph{ampiezza} (la differenza fra i due stati come vettori) scala
come $1/N$; l'infedeltà è il \emph{modulo quadro} di quella differenza,
quindi scala come $(1/N)^2=1/N^2$ -- raddoppiando $N$, l'infedeltà si
riduce di un fattore $2^2=4$, non $2$. Il fatto che il rapporto misurato
converga proprio a $4$ (non a $3$, non a $5$) è la prima conferma
numerica, indipendente dalla teoria, che l'implementazione è
effettivamente un metodo di Trotter del prim'ordine.

\subsection{Self-test 4: fattorizzazione esatta di $H_0$}

L'ultimo test verifica che $H_0=b\sum_iZ_i+H_{ex}$ si fattorizzi
\emph{esattamente} come campo $\times$ scambio, indipendentemente da
$\tau$: per tre set di parametri casuali, $\|e^{-iH_0\tau} -
e^{-ib\tau S_z^{tot}}e^{-iH_{ex}\tau}\|\sim10^{-16}$. La ragione teorica
(perché questo blocco non ha \emph{nessun} errore di approssimazione, a
differenza di tutto il resto) è derivata per esteso nella
Sezione~\ref{sec:trotter-error}.

\section{Ricerca del punto di lavoro (Sezione 2 del notebook)}
\label{sec:ricerca-punto}

\subsection{Perché serve il termine DM: verifica diretta}

Prima ancora di cercare un buon punto, va stabilito che il termine DM è
\emph{necessario} per avere una qualunque dinamica osservabile. Si
verifica esplicitamente che $\ket{000}$ è autostato esatto di $H_{ex}$:
\begin{equation}
H_{ex}\ket{000} = (J+2J')\,\ket{000},
\label{eq:000-eigenstate}
\end{equation}
un fatto che si dimostra così: il termine di scambio $X_iX_j+Y_iY_j$ fra
due spin entrambi "su" (come sono tutti e tre in $\ket{000}$) si annulla
identicamente. Questo perché $X_iX_j+Y_iY_j = 2(\sigma_i^+\sigma_j^- +
\sigma_i^-\sigma_j^+)$ (operatori di innalzamento/abbassamento), e
applicato a $\ket{\uparrow\uparrow}$: il termine $\sigma_i^+\sigma_j^-$
abbassa il sito $j$ (già su, quindi produce $\ket{\uparrow\downarrow}$)
ma poi il fattore $\sigma_i^+$ tenta di alzare il sito $i$ che è già su
-- risultato nullo (non si può alzare oltre $\ket\uparrow$); simmetricamente
per l'altro termine. Resta solo $Z_iZ_j\ket{\uparrow\uparrow}=+1\cdot\ket{\uparrow\uparrow}$
per ciascuna coppia, da cui l'Eq.~\eqref{eq:000-eigenstate} (verificato
numericamente nel documento di supporto, errore $2.2\times10^{-16}$).
Questo identifica $\ket{000}$ come il membro di punta ($M=3/2$) del
multipletto $A$ della decomposizione di Kambe (\texttt{\seqsplit{trimer\_ring\_exact.py}}),
con $E_A=J+2J'$.

Aggiungendo il campo, $\ket{000}$ resta autostato anche di
$H_0=b\sum_iZ_i+H_{ex}$ (perché $S_z^{tot}\ket{000}=3\ket{000}$ commuta
banalmente con l'essere già autostato di $H_{ex}$), con autovalore
$3b+J+2J'$ -- verificato, errore $4.4\times10^{-16}$. \textbf{Conseguenza
diretta}: se $D=0$, allora $H=H_0$ e $\ket{000}$ è autostato esatto
dell'Hamiltoniana completa, quindi $\ket{\psi(t)}=e^{-iEt}\ket{000}$
(solo una fase globale) e $\langle S_z^{tot}\rangle(t)$ è \textbf{costante
nel tempo}, per qualunque $b,J,J'$ -- verificato numericamente su una
griglia di tempi, scarto $0$ a precisione macchina. Senza il termine DM
non c'è alcuna dinamica da osservare: è $D\neq0$, e solo quello, a rendere
possibile un'oscillazione. Lo stesso fenomeno, con la stessa dimostrazione
concettuale, era già stato stabilito per il dimero
($\ket{00}=\ket{t_+}$ autostato di $H_1$).

\subsection{Il criterio a due indicatori}

Fissati $J=1,J'=0.4$ (punto di lavoro del VQE), bisogna scegliere $(b,D)$
in modo che $\langle S_z^{tot}\rangle(t)$ mostri oscillazioni
\textbf{non banali}: non un singolo coseno pulito (che metterebbe alla
prova a malapena il codice, essendo equivalente a un sistema a due
livelli), ma la sovrapposizione di più frequenze indipendenti.

\subsubsection{Decomposizione spettrale (formula usata da \texttt{mode\_amplitudes})}

Sia $H\ket{k}=E_k\ket{k}$ la base di autostati esatti (ottenuta nel
notebook via \texttt{numpy.linalg.eigh}), e $c_k=\langle k\ket{000}$ le
proiezioni dello stato iniziale su questa base. Lo stato evoluto è
$\ket{\psi(t)}=\sum_k c_k e^{-iE_kt}\ket k$, e il valore di aspettazione
di un qualunque operatore hermitiano $\hat O$ (qui $\hat O=S_z^{tot}$) è
\begin{align}
\langle\hat O\rangle(t) &= \sum_{k,l} c_k^*c_l\,e^{i(E_k-E_l)t}\,O_{kl}
= \sum_k |c_k|^2 O_{kk} + \sum_{k\neq l} c_k^*c_l\,e^{-i(E_l-E_k)t}\,O_{kl}\\
&= \underbrace{\sum_k|c_k|^2O_{kk}}_{\text{parte costante}} + \sum_{k<l} 2\,\mathrm{Re}\!\left[c_k^*c_l\,O_{kl}\,e^{-i\Delta_{kl}t}\right], \qquad \Delta_{kl}\equiv E_l-E_k,
\label{eq:mode-decomp}
\end{align}
dove nell'ultimo passaggio si sono accoppiati i termini $(k,l)$ e
$(l,k)$ usando $O_{lk}=O_{kl}^*$ ($\hat O$ hermitiano) e
$c_l^*c_k\,e^{i\Delta_{kl}t}=(c_k^*c_l\,e^{-i\Delta_{kl}t})^*$. L'Eq.~\eqref{eq:mode-decomp}
è esattamente la funzione \texttt{\seqsplit{mode\_amplitudes}} del notebook: per
ogni coppia $k<l$ calcola l'\textbf{ampiezza del modo}
$a_{kl}=2|c_k^*c_l\,O_{kl}|$ e il suo \textbf{gap} $\Delta_{kl}$, ordinando
per ampiezza decrescente. I due modi con ampiezza più alta, $a_1\ge a_2$,
sono quelli che dominano la forma di $\langle S_z^{tot}\rangle(t)$.

\subsubsection{I due indicatori e il loro significato fisico}

\begin{enumerate}
\item \textbf{Escursione picco-picco} di $\langle S_z^{tot}\rangle(t)$ su
una finestra temporale lunga (il notebook usa $t_{\max}=30$): il segnale
deve essere abbastanza grande da restare leggibile sopra il rumore
statistico di un campionamento a shot finiti, che scala come
$1/\sqrt{N_{\rm shots}}$ (per $N_{\rm shots}=8192$, soglia
$\approx0.011$).
\item \textbf{Rapporto $a_2/a_1$}: se $a_2/a_1\to0$, un solo modo domina
e il segnale è un coseno quasi puro (\emph{monocromatico}, nel senso
ottico del termine: una sola frequenza); se $a_2/a_1\to1$, due modi con
ampiezza comparabile producono \textbf{battimenti} -- massimi e minimi
irregolari, senza periodicità semplice, la firma di una dinamica
autenticamente multi-livello.
\end{enumerate}
Servono \textbf{entrambi}: un segnale grande ma monocromatico (es.\ vicino
a una risonanza a due livelli) non soddisferebbe il secondo criterio; un
segnale spettralmente ricco ma minuscolo sarebbe inutilizzabile
sperimentalmente. Lo stesso criterio, con la stessa motivazione, è stato
usato per derivare i punti R0/R1 del dimero.

\subsection{Lo scan e i candidati trovati}

Il notebook scandisce una griglia $60\times40$ in $(b,D)$, con
$b\in[0.05,5]$, $D\in[0.05,2.5]$ (2400 punti), calcolando per ciascuno
picco-picco e $a_2/a_1$ dalla dinamica \textbf{esatta} (diagonalizzazione,
non Trotter -- a questo stadio si sta solo scegliendo il punto di lavoro,
non validando l'algoritmo di Trotter). I primi candidati, filtrati per
picco-picco $>0.5$ e ordinati per $a_2/a_1$ decrescente:

\begin{center}
\begin{tabular}{@{}ccccc@{}}
\toprule
$b$ & $D$ & $D/J$ & picco-picco & $a_2/a_1$ \\
\midrule
0.05 & 1.93 & 1.93 & 3.80 & 0.9998 \\
0.72 & 1.24 & 1.24 & 0.51 & 0.9992 \\
0.30 & 1.31 & 1.31 & 1.38 & 0.9982 \\
\bottomrule
\end{tabular}
\end{center}

\subsection{Verifica di robustezza: perché non basta il primo posto in classifica}

Un candidato preso senza altre verifiche potrebbe essere un punto
\emph{fine-tuned} (un picco stretto e isolato nel piano $(b,D)$,
privo di significato fisico robusto) piuttosto che un vero
\emph{plateau} fisico. Il notebook perturba ciascun candidato di
$\pm5$--$10\%$ in $b$ e $D$ e ricalcola i due indicatori:

\begin{center}
\begin{tabular}{@{}cc@{}}
\toprule
punto & comportamento sotto perturbazione \\
\midrule
$(0.05,1.93)$ & picco-picco resta $3.7$--$3.8$; $a_2/a_1$ scende a $0.74$--$0.91$ \\
$(0.30,1.31)$ & picco-picco $1.1$--$1.8$; $a_2/a_1$ $0.75$--$0.99$ \\
$(0.72,1.24)$ & picco-picco $0.4$--$0.65$; $a_2/a_1$ $0.6$--$0.98$ \\
\bottomrule
\end{tabular}
\end{center}
Tutti e tre restano qualitativamente validi (nessun collasso improvviso
del segnale): sono plateau, non punte isolate. Questo giustifica
l'adozione del primo, che ha anche il segnale assoluto più grande.

\section{Il punto adottato: $R_0$ (trimero)}

$$J=1, \qquad J'=0.4, \qquad b=0.05, \qquad D=1.93.$$
\textbf{Da dichiarare esplicitamente al relatore come provvisorio.} Per
il dimero, il punto analogo ($R_1$) era stato derivato dalla struttura a
V dei quattro livelli (formula di Rabi, due detuning indipendenti), non
da uno scan. Per il trimero la stessa derivazione richiederebbe
classificare le transizioni fra gli 8 livelli di Kambe sotto la rottura
di $S_{12}^2$ indotta dall'Opzione B del DM -- un lavoro più lungo,
pianificato ma non ancora fatto. Il punto $R_0$ (trimero) qui adottato è
uno scan \emph{verificato robusto}, non un semplice primo posto in una
classifica cieca (si veda la nota procedurale nel documento di
applicazione sulla stessa lezione, già imparata una volta con un errore
simile nel dimero e qui non ripetuta).

\section{Convergenza in $N$ (Sezione 4 del notebook)}
\label{sec:trotter-error}

Questa è la parte concettualmente più densa del notebook: va spiegato
\emph{da dove viene} l'errore di Trotter, non solo mostrarne i numeri.

\subsection{Perché il passo di Trotter ha tre livelli, non uno}

Il passo implementato è
\begin{align}
U_{\rm step}(\tau) &= V_{DM}(\tau)\;e^{-ib\tau S_z^{tot}}\;V_{ex}(\tau), \label{eq:Ustep}\\
V_{ex}(\tau)&=U_{31}(\tau)\,U_{23}(\tau)\,U_{12}(\tau), \qquad V_{DM}(\tau)=U_{31}^{DM}(\tau)\,U_{23}^{DM}(\tau)\,U_{12}^{DM}(\tau), \notag
\end{align}
dove $U_{ij}(\tau)=e^{-iJ_{ij}\tau\vec\sigma_i\cdot\vec\sigma_j}$ e
$U_{ij}^{DM}(\tau)$ sono i blocchi di \emph{singolo} legame, esatti (Sezione~2
di \texttt{\seqsplit{quantum\_simulation\_trimero\_anello\_teoria.tex}}: la
dimostrazione usa che $X_iX_j,Y_iY_j,Z_iZ_j$ commutano a due a due per una
coppia di qubit fissata, quindi il loro esponenziale fattorizza senza
errore). L'evoluzione esatta è invece $U_{\rm esatto}(\tau)=e^{-i(H_0+H_{DM})\tau}$.

Perché $U_{\rm step}\neq U_{\rm esatto}$? Perché la costruzione
dell'Eq.~\eqref{eq:Ustep} introduce \textbf{tre} approssimazioni
indipendenti, non una sola:
\begin{enumerate}
\item[\textbf{Livello 0}] (esterno): si tratta $H_0$ e $H_{DM}$ come due
blocchi separati in sequenza, $e^{-i(H_0+H_{DM})\tau}\approx
e^{-iH_{DM}\tau}e^{-iH_0\tau}$, il che è esatto solo se $[H_0,H_{DM}]=0$
-- non è il caso (il campo non commuta con il DM, che rompe la
conservazione di $S_z^{tot}$).
\item[\textbf{Livello 1}] (interno, scambio): $H_{ex}$ è la somma di
\emph{tre} legami che condividono i qubit a due a due (il qubit 2 sta sia
nel legame $(1,2)$ sia nel legame $(2,3)$); si dimostra
(\texttt{\seqsplit{quantum\_simulation\_trimero\_anello\_teoria.tex}} \S4)
che $[\vec\sigma_1\!\cdot\!\vec\sigma_2,\vec\sigma_2\!\cdot\!\vec\sigma_3]\neq0$,
quindi $V_{ex}(\tau)\neq e^{-iH_{ex}\tau}$ esattamente.
\item[\textbf{Livello 2}] (interno, DM): per lo stesso motivo (qubit
condivisi), anche i tre legami DM non commutano fra loro
(\texttt{\seqsplit{quantum\_simulation\_trimero\_anello\_teoria.tex}} \S6), quindi
$V_{DM}(\tau)\neq e^{-iH_{DM}\tau}$ esattamente.
\end{enumerate}
Il livello 2 è la scoperta di questo lavoro rispetto al compito
originale (che prevedeva solo il livello 1): non era affatto scontato che
anche $H_{DM}$ avesse lo stesso problema strutturale dello scambio, ed è
stato verificato esplicitamente prima di essere assunto.

\subsection{Da dove viene l'errore: la formula di Baker--Campbell--Hausdorff}

Per due operatori $A,B$ generici, non necessariamente commutanti,
\begin{equation}
e^Ae^B = \exp\!\Big(A+B+\tfrac12[A,B]+\tfrac1{12}[A,[A,B]]+\tfrac1{12}[B,[B,A]]+\cdots\Big).
\label{eq:BCH}
\end{equation}
Se $[A,B]=0$, ogni termine della serie oltre il primo si annulla (per
induzione: ogni commutatore annidato successivo è costruito a partire da
$[A,B]$ o da altri commutatori già nulli), e $e^Ae^B=e^{A+B}$
\emph{esattamente}. È questo il motivo per cui i blocchi di singolo
legame (Sezione~\ref{subsec:selftest2}) e la fattorizzazione di $H_0$
(self-test 4) sono esatti: in entrambi i casi il commutatore rilevante è
identicamente nullo. Se invece $[A,B]\neq0$, approssimare
$e^{A+B}\approx e^Ae^B$ significa scartare il termine $\tfrac12[A,B]$ e
tutti i successivi.

\subsection{L'errore per singolo passo, in formula}

Espandendo entrambi i membri dell'Eq.~\eqref{eq:BCH} in serie di Taylor
fino a $O(\tau^2)$ (per $A\to A\tau$, $B\to B\tau$, $\tau$ piccolo):
\begin{equation}
e^{-i(A+B)\tau} - e^{-iA\tau}e^{-iB\tau} = -\frac{\tau^2}2[A,B] + O(\tau^3)
\label{eq:2term-error}
\end{equation}
(derivazione per esteso: si espandono $e^{-iA\tau}$, $e^{-iB\tau}$ e
$e^{-i(A+B)\tau}$ ciascuno fino al second'ordine in $\tau$, si moltiplicano
i primi due, e si confronta termine a termine col terzo -- fatto per
esteso in \texttt{\seqsplit{quantum\_simulation\_dimero\_teoria.tex}} \S6.1 per il
caso a due fattori, e generalizzato a tre fattori nella teoria del
trimero). L'errore per passo è quindi $O(\tau^2)$, \emph{proporzionale al
commutatore}: se $[A,B]=0$ l'errore si annulla identicamente, qualunque
sia $\tau$ -- coerente con quanto osservato nei self-test.

Per il caso del trimero, applicando questa formula ai tre livelli
sopra descritti e sommando i contributi per disuguaglianza triangolare
(la derivazione completa, incluso il calcolo esplicito dei commutatori
tramite l'operatore di chiralità di spin
$\chi=\vec\sigma_1\cdot(\vec\sigma_2\times\vec\sigma_3)$, è in
\texttt{\seqsplit{quantum\_simulation\_trimero\_anello\_teoria.tex}} \S8), si ottiene
che l'errore per passo resta $O(\tau^2)$ nonostante i tre contributi
indipendenti:
\begin{equation}
\big\|U_{\rm esatto}(\tau)-U_{\rm step}(\tau)\big\| \le \|\varepsilon_0(\tau)\|+\|\varepsilon_1(\tau)\|+\|\varepsilon_2(\tau)\| = O(\tau^2).
\label{eq:total-step-error}
\end{equation}

\subsection{Dall'errore per passo all'errore totale: perché il rapporto tende a 4}

Per $N$ passi di durata $\tau=t/N$, usando la disuguaglianza triangolare
$N$ volte sull'Eq.~\eqref{eq:total-step-error}:
\begin{equation}
\big\|U_{\rm esatto}(t)-U_{\rm step}(\tau)^N\big\| \lesssim N\cdot O(\tau^2) = N\cdot O\!\left(\frac{t^2}{N^2}\right) = O\!\left(\frac{t^2}{N}\right).
\label{eq:total-N-error}
\end{equation}
L'errore sull'\emph{operatore} (e quindi, a meno di dettagli, sull'ampiezza
dello stato) scala come $1/N$. L'\textbf{infedeltà} misurata nel notebook,
$1-|\langle\psi_{\rm esatto}|\psi_{\rm Trotter}\rangle|^2$, è legata al
\emph{quadrato} della differenza fra i due stati (uno sviluppo di Taylor
di $|\langle\psi_1|\psi_2\rangle|^2$ attorno a $\psi_1=\psi_2$ dà
$1-|\langle\psi_1|\psi_2\rangle|^2\approx\|\psi_1-\psi_2\|^2$ al
prim'ordine), quindi scala come $(1/N)^2=1/N^2$: raddoppiando $N$,
l'infedeltà scende di un fattore $4$. È esattamente il rapporto misurato
nel self-test 3 e nella tabella di convergenza principale.

\subsection{I numeri del notebook, con questa lettura}

\begin{center}
\begin{tabular}{@{}ccc@{}}
\toprule
$N$ & infedeltà finale (a $T\simeq7.85$, 2 periodi) & rapporto raddoppiando $N$ \\
\midrule
80 & $1.665\times10^{-1}$ & \\
160 & $4.011\times10^{-2}$ & $4.15$ \\
320 & $1.007\times10^{-2}$ & $3.98$ \\
640 & $2.545\times10^{-3}$ & $3.96$ \\
1280 & $6.409\times10^{-4}$ & $3.97$ \\
\bottomrule
\end{tabular}
\end{center}
Il grafico log--log del notebook mostra questi punti allineati con la
retta di riferimento $\propto1/N^2$, la firma grafica diretta
dell'Eq.~\eqref{eq:total-N-error} elevata al quadrato. Il notebook
riporta anche gli $N$ minimi per soglie di infedeltà standard: $<1\%$ a
$N\simeq325$, $<0.1\%$ a $N\simeq1025$, $<0.01\%$ a $N\simeq3250$ --
sensibilmente più alti di quanto bastasse per il dimero ($N=80$ dava già
$\sim10^{-5}$), coerentemente con la presenza di due livelli di errore in
più (1 e 2) rispetto al solo livello 0 del dimero, e con $D/J\simeq1.9$
non piccolo (l'errore di livello 1 scala con $J'^2$, quello di livello 2
con $D^2\propto r^2$: entrambi crescono con l'intensità delle interazioni
coinvolte).

\section{Separazione dei contributi di errore (Sezione 5 del notebook)}

Il relatore ha chiesto esplicitamente, oltre alla convergenza in $N$, di
separare -- se possibile -- i contributi di errore. Poiché il livello 0 e
il livello 1 erano già presenti nel disegno originale del compito (Trotter
esterno fra campo/scambio e DM, Trotter interno per lo scambio), mentre il
livello 2 (Trotter interno per il DM) è una scoperta di questo lavoro, il
confronto più informativo è isolare \emph{quanto pesa} il livello 2 da
solo.

\subsection{Il confronto $V_1$ vs $V_2$}

Si costruiscono due varianti del passo:
\begin{itemize}
\item $V_2$ (\textbf{reale}, quella effettivamente implementata nel
circuito): $H_{DM}$ spezzato nei tre legami, come in Eq.~\eqref{eq:Ustep};
\item $V_1$ (\textbf{idealizzata}, non realizzabile con soli gate
fisicamente etichettati): $H_{DM}$ trattato come un singolo blocco esatto,
$e^{-iH_{DM}\tau}$, calcolato con \texttt{scipy.linalg.expm} sull'intera
matrice $8\times8$ -- possibile solo perché il sistema è ancora piccolo
abbastanza da diagonalizzare classicamente, non un'opzione per sistemi
grandi (si veda \texttt{\seqsplit{trimero\_anello\_circuito\_compatto.tex}} per
la discussione completa di perché questa scorciatoia non generalizza).
\end{itemize}

\begin{center}
\begin{tabular}{@{}cccc@{}}
\toprule
$N$ & infedeltà $V_2$ (reale) & infedeltà $V_1$ (idealizzata) & rapporto \\
\midrule
320 & $1.007\times10^{-2}$ & $2.306\times10^{-3}$ & $4.370$ \\
640 & $2.545\times10^{-3}$ & $5.788\times10^{-4}$ & $4.397$ \\
1280 & $6.409\times10^{-4}$ & $1.452\times10^{-4}$ & $4.413$ \\
\bottomrule
\end{tabular}
\end{center}

Il rapporto si stabilizza a $\simeq4.4$: il livello 2, se rimosso
(idealmente), ridurrebbe l'infedeltà totale di un fattore $\sim4.4$ a
questo punto di lavoro. Non è un contributo trascurabile.

\subsection{Un avvertimento onesto sul segno}

La teoria (Sezione~\ref{sec:trotter-error}) garantisce solo che ciascun
livello sia $O(\tau^2)$ in \emph{modulo} (disuguaglianza triangolare, un
limite superiore); non garantisce che i tre contributi si sommino sempre
costruttivamente. Verificato esplicitamente (non assunto) in un'analisi
collaterale: cambiando punto di lavoro e ordine dei fattori nel prodotto
di Trotter, il rapporto $V_2/V_1$ misurato varia da $0.62$ a $1.57$ --
talvolta il livello 2 riduce lievemente l'errore totale (cancellazione
parziale fra i termini di commutatore), talvolta lo aumenta. A $R_0$
(trimero) il rapporto è $\simeq4.4$, ma questo valore specifico non va
generalizzato ad altri punti di lavoro senza rifare il confronto.

\section{Il circuito (Sezione 6 del notebook)}

Il notebook disegna il circuito per 2 passi completi (separati da un
barrier), nell'ordine $H_{ex}$ (3 legami) $\to$ campo (esatto, 3 rotazioni
$R_z$) $\to$ $H_{DM}$ (3 legami). Conteggio per un singolo passo: 12
Hadamard, 9 $R_{ZZ}$, 3 $R_{XX}$, 3 $R_{YY}$, 3 $R_z$ -- 15 gate a due
qubit nativi. Ai valori di $N$ trovati per la convergenza ($\sim300$ a
$3000$), il circuito totale avrebbe da $4500$ a $45000$ gate a due qubit:
irrilevante per la simulazione statevector di adesso (puro calcolo
classico), ma il numero da tenere a mente quando si affronterà la Parte 2
(rumore) -- è lo stesso tipo di conteggio che rende rilevante, in quella
sede, il vantaggio di gate di RBS su $W$ già documentato per il VQE.

\section{Conclusioni}

\begin{enumerate}
\item Il modulo è validato a precisione macchina su quattro test
indipendenti, incluso il confronto diretto fra il circuito Qiskit reale e
la costruzione a matrice.
\item La struttura dell'errore di Trotter ha \textbf{tre} livelli
indipendenti (non due, come inizialmente previsto dal compito): il campo
non commuta col DM (livello 0), i tre legami di scambio non commutano fra
loro (livello 1, noto), e i tre legami DM neppure (livello 2, scoperta di
questo lavoro) -- tutti e tre derivati in forma chiusa, non solo verificati
a numeri.
\item Il punto di lavoro $R_0$ (trimero) usato per la dimostrazione è
\textbf{provvisorio}: un plateau verificato robusto, non (ancora) una
derivazione fisica principiata come quella fatta per il dimero.
\item La convergenza misurata ($N\sim300$--$3000$ per soglie di infedeltà
dall'1\% allo 0.01\%) e il peso non trascurabile del livello 2
($\sim4.4\times$) sono i due numeri principali da portare al relatore in
risposta diretta alla sua richiesta di "testare quanti passi $N$ servono"
e di "separare, se possibile, i contributi di errore".
\end{enumerate}

\end{document}
